更多细节见「同上」。此猜想纯属组合，但 Lusztig 与 Yun \cite{LY} 的证明是几何的，并且实际上用到了等特征几何 Satake！我们于是走向相反的方向：证明该公式蕴含上述交换性约束是正确的。

因此我们在混合特征情形对定理 \ref{intro:geomSat} 的证明，用到了等特征几何 Satake。一个有趣的问题是直接证明上述组合公式，那将给出几何 Satake 在等特征与混合特征情形下的纯局部证明。此外，类似的策略应当也适用于证明分歧 $p$ 进群的几何 Satake（其等特征对应版本已在 \cite{Zh0, Ri} 中建立）。


\subsection{动机 Satake 范畴}
\label{SS: MotSat}
为解释几何 Satake 应用于志村簇算术几何背后的某些想法，我们想用动机的语言重新表述几何 Satake。
本小节须假定读者熟悉（纯）动机的基本理论。经典参考是 \cite{Kl}。

本小节假定 $k$ 是完全域。
我们将构造（数值）动机范畴中一个 $\bQ$-线性、半单、超 Tannakian 的子范畴 $\Sat^m$。把 $\Sat^m$ 与 $\Ql$ 作张量积，可以恢复（通常 Satake 范畴的一个变体）。
为避免某些复杂的几何，我们主要关注 $\GL_n$ 情形，并在本小节末尾简述处理一般约化群所需的修改。



令 $\on{Var}_k$ 表示 $k$ 上光滑射影簇的范畴。
令 $\on{Mot}^{\on{rat}}_k$（相应地 $\on{Mot}_k^{\on{num}}$）表示 $k$ 上关于有理（相应地数值）等价关系的纯动机构成的 $\bQ$-线性范畴。则 $\on{Mot}_k^*$（$*=\on{rat},\on{num}$）是对称幺半的 $\bQ$-线性伪 Abel 范畴，而 $\on{Mot}_k^{\on{num}}$ 是超 Tannakian 半单 $\bQ$-线性 Abel 范畴（见 \cite{Ja}）。依 \cite{Ja}，$\on{Mot}_k^{*}$ 中的对象记作三元组 $(X,p,m)$，其中 $X$ 是 $k$ 上光滑射影簇，$p$ 是 $X$ 的零次自对应环中的幂等元（用上述某个适当等价关系定义），$m\in\bZ$。于是有自然函子
\[h:\on{Var}_k^{\on{op}}\to \on{Mot}^*_k,\quad X\mapsto h(X):=(X,\id,0).\]
照常，把 $(X,\id,m)$ 写作 $h(X)(m)$。
函子 $h$ 经 $\on{Var}_k$ 关于普遍同胚的局部化分解。
注意 $\cha k=p>0$，从 $\on{Var}_k$ 到完全 $k$-概形范畴 $\on{Sch}_k^\pf$ 的完全化函子经此局部化分解，而且实际上诱导该局部化范畴到 $\on{Sch}_k^\pf$ 的全嵌入（由 \cite[推论 A.16、命题 A.17]{Zh1}，另见 \cite[引理 3.8、命题 3.11]{BS}）。其本质像记作 $\on{Var}^\pf_k$，由作为光滑射影概形之完全化的 $k$ 上完全概形组成。
因此得到函子
$$h:(\on{Var}^\pf_k)^{\on{op}}\to\on{Mot}^*_k.$$

回忆给定序列 $\mmu=(\mu_1,\ldots,\mu_r)$，有 Schubert 簇的扭积 $\Gr_{\mmu}=\Gr_{\mu_1}\tilde\times\cdots\tilde\times\Gr_{\mu_r}$，其维数记作 $d_\mmu$。当且仅当每个 $\mu_i$ 极小，$\Gr_\mmu$ 才是某光滑射影簇的完全化，从而定义 $\on{Var}^\pf_k$ 中的对象。在 $\GL_n$ 情形，$\Gr_\mmu$ 是逐次以完全 Grassmann 簇为纤维的纤维化，
分类满足
$$\La_0=\mO^n, \quad  \varpi^{d_i+1}\La_{i-1}\subset \La_i\subset \varpi^{d_i}\La_{i-1} \mbox{ （对某个 } d_i\in\bZ\mbox{）}$$
的格链 $\{\La_i, i=0,\ldots, r\}$。
（这里及本小节余下部分，为简便只描述 $\Gr_\mmu$ 的 $k$-点。）
令 $\on{H}(\Gr_\mmu, \Gr_{\nu_\bullet})$ 为以 $\Gr_{\mmu}\times_\Gr\Gr_{\nu_\bullet}$ 的、维数\textbf{恰为} $\frac{1}{2}(d_{\mmu}+d_{\nu_\bullet})$ 的不可约分支为基的 $\bQ$-向量空间。我们需要如下几何事实。
\begin{lem}
\label{L:faithful realization}
设 $m,n,r$ 为三个整数，满足 $2(m-n)=d_{\mmu}-d_{\nu_\bullet}$ 与 $2(n-r)=d_{\nu_\bullet}-d_{\la_\bullet}$。
则自然映射
\[\on{H}(\Gr_{\mmu},\Gr_{\nu_\bullet})\to \Hom_{\on{Mot}^*_k}(h(\Gr_{\mmu})(m),h(\Gr_{\nu_\bullet})(n))\]
是单射。此外，存在唯一的 $\bQ$-线性映射
$$\on{H}(\Gr_{\mmu},\Gr_{\nu_\bullet})\otimes \on{H}(\Gr_{\nu_\bullet},\Gr_{\la_\bullet})\to \on{H}(\Gr_{\mmu},\Gr_{\la_\bullet})$$ 
使下图交换
{\scriptsize
\[
\begin{tikzcd}[row sep=2.2em, column sep=0.3em, thick]
\on{H}(\Gr_{\mmu},\Gr_{\nu_\bullet})\otimes \on{H}(\Gr_{\nu_\bullet},\Gr_{\la_\bullet}) \arrow[r] \arrow[d] & \on{H}(\Gr_{\mmu},\Gr_{\la_\bullet}) \arrow[d] \\
\Hom_{\on{Mot}^*_k}(h(\Gr_{\mmu})(m),h(\Gr_{\nu_\bullet})(n))\otimes\Hom_{\on{Mot}^*_k}(h(\Gr_{\nu_\bullet})(n),h(\Gr_{\la_\bullet})(r)) \arrow[r] & \Hom_{\on{Mot}^*_k}(h(\Gr_{\mmu})(m),h(\Gr_{\la_\bullet})(r))
\end{tikzcd}
\]
}
\end{lem} 
\begin{proof}
我们只给出梗概。
由于 $\Gr_\mmu\times\Gr_{\nu_\bullet}$ 允许胞腔分解（用半无限轨道的卷积积，例如见 \cite[\S 3.2.5, \S 3.2.6]{XZ}），此时有理等价、同调等价与数值等价三者一致。于是只需证明 $\Gr_{\mmu}\times_\Gr\Gr_{\nu_\bullet}$ 的这些 $\frac{1}{2}(d_{\mmu}+d_{\nu_\bullet})$ 维不可约分支的轮类，在 $\Gr_{\mmu}\times\Gr_{\nu_\bullet}$ 的 étale 上同调中线性无关。而这可由卷积映射 \eqref{conv m2}（它是半小映射）的通常分解定理推出。

给定单射性以及对应对合成的定义，要证图的交换性，只需注意在 $\Gr_{\mmu}\times\Gr_{\nu_\bullet}\times\Gr_{\la_\bullet}$ 中，来自 $\on{H}(\Gr_{\mmu},\Gr_{\nu_\bullet})$ 与 $\on{H}(\Gr_{\nu_\bullet},\Gr_{\la_\bullet})$ 的轮之相交是真交，然后应用维数估计 \eqref{E:dim est}。
\end{proof}



现在定义 $\bQ$-线性加性范畴 $\Sat_G^{0}$，并把\textbf{动机 Satake 范畴} $\Sat_G^m$ 定义为 $\Sat_G^{0}$ 的幂等完备化。我们指出，定义 $\Sat_G^0$ 的想法已见于 \cite[\S 4.3]{FKK}。
\begin{itemize}
\item 对象是数对 $(X,f)$，其中 $X$ 是形如 $\Gr_\mmu$ 的完全光滑射影概形的不交并（$\mmu$ 取遍极小余权序列），$f$ 是 $X$ 上取值于 $\bZ$ 的局部常值函数。若 $X=\Gr_\mmu$，则 $f$ 由一个整数 $m$ 给出；此时把该数对记作 $(\Gr_\mmu,m)$。若还有 $m=0$，有时简记 $(\Gr_\mmu,0)$ 为 $\Gr_\mmu$。
\item \[\Hom_{\Sat_G^0}((\Gr_{\mmu},m),(\Gr_{\nu_\bullet},n))=\left\{\begin{array}{ll}0 & 2(m-n)\neq d_\mmu-d_{\nu_\bullet}\\
\on{H}(\Gr_{\mmu},\Gr_{\nu_\bullet}) & 2(m-n)=d_\mmu-d_{\nu_\bullet} \end{array}\right.\]
若把 $(X,f)=\sqcup_i (X_i,m_i)$ 与 $(Y,g)=\sqcup (Y_j,n_j)$ 写成连通对象的不交并，则定义
$$\Hom_{\Sat_G^0}((X,f),(Y,g))=\bigoplus_{ij}\Hom_{\Sat_G^0}((X_i,m_i),(Y_j,n_j)).$$
\end{itemize}
